\chapter{Statustafel}

Rang der Befunde. \emph{Belegt} heißt: an den herangezogenen Texten wörtlich ausgewiesen. \emph{Erschlossen} heißt: aus belegten Sätzen mehrerer Autoren zusammengeführt; die Zusammenführung ist eigene Deutung. \emph{Offen} heißt: nicht entschieden, absichtlich.

\begin{itemize}
\item Belegt: Das zeitliche Kontinuum besteht nur als Grenze, das räumliche als Ganzes (Brentano).
\item Belegt: In der Zeit ist ein Punkt der Sache nach ausgezeichnet, im Raum keiner (Brentano, S.\,117).
\item Belegt: Der Raum lässt sich überall begrenzen, aber nicht unterbrechen (Hegel, Nachschrift).
\item Belegt: Die Zeit ist der wahrhafte Punkt; der Ort ist das räumliche Jetzt (Hegel, Nachschrift und Enzyklopädie).
\item Belegt: Die Zeit ist die Faserrichtung der Kausalketten, der Raum das Nebeneinander der Faserbündel (Reichenbach, \S\,43).
\item Belegt: Die Zeit ist gegenüber dem Raum das logisch Primäre; das Lichtjahr ist die logische Urform aller Längenmessung (Reichenbach, \S\,27, \S\,42).
\item Belegt: Die Dimensionszahl des Raumes ist der objektive Sinn des Nahwirkungsprinzips (Reichenbach, \S\,44).
\item Belegt: Die Einsinnigkeit der Zeit bleibt bei Reichenbach ungelöst (\S\,21).
\item Belegt: Räumliche Differenzen sind sachlich, zeitliche modal (Brentano, VII--VIII); später abgeschwächt (Brentano, X).
\item Belegt: Das Jetzt ist weltweit, das Hier je meines (Fink).
\item Belegt: Wir selbst sind die notwendige Garantie der indexikalischen Spezifizierbarkeit (Koch 1994).
\item Belegt: Die Richtungen des Raumes kommen vom Leib, die Modi der Zeit von den Seelenvermögen (Koch 2010).
\item Belegt: Die Zeit wird zum Raum, nicht umgekehrt; die Zeit ist immer schon eingeräumt (Guzzoni, mit Heidegger).
\item Belegt: Leibniz' Definition des Raumes setzt das Räumliche voraus (Baumann).
\item Belegt: Die Zeitvorstellung wird auf das dauernde Ich bezogen, dessen Dauer der Zeit vorausgeht (Baumann, S.\,659--668).
\item Belegt: Die Zahl ist weder aus der Zeit noch aus dem Raum abzuleiten (Baumann, S.\,668--669).
\item Belegt: Der Raum setzt die Zeit voraus und ist nur als Zeitraum konzipierbar; die Zeit wird nur am Raum zur Größe (Martić, zu Kant).
\item Belegt: Um die Zukunft zu erfahren, muss man nur warten (Chiba).
\item Belegt: Das negative Vorzeichen ist Ausdruck der Sonderstellung der Zeit, nicht deren Aufhebung (Reichenbach; Schneider zu Minkowski).
\item Erschlossen: Der Chiasmus der Grenze, ein ausgezeichneter Schnitt ohne beliebige gegen beliebige Schnitte ohne ausgezeichneten; Fink und Baumann widersprechen einander darin nicht.
\item Erschlossen: Das weltweite Jetzt ist die Zeit in der Gestalt des Raumes und als solches definitorisch; das nicht-definitorische Jetzt ist das an der Faser.
\item Erschlossen: Die Richtungen des Raumes, bis zur Dimensionszahl, stammen aus der Wirkung, also aus dem Zeitlichen.
\item Erschlossen: Die Abstraktion vom Standpunkt ist die Abstraktion von der Zeit; der zeit-transzendente Standpunkt macht die Zeit zum Raum.
\item Erschlossen: Baumanns Ich-Dauer, Reichenbachs ewiges Jetzt und Hegels Ich = Ich sagen dasselbe.
\item Erschlossen: Hegels Übergang Raum--Zeit und Guzzonis Übergang Zeit--Raum widersprechen sich nicht; beide ergeben dieselbe Rangordnung.
\item Erschlossen: Die Signatur ist die metrische Spur der Asymmetrie, nicht die Asymmetrie.
\item Offen: Was der Raum ist, sofern er nicht Grenze, Richtung und Maß hat; ob er Rest der Ableitung oder eigener Art ist.
\item Offen: Die Einsinnigkeit der Zeit; warum die Zukunft von selbst kommt.
\item Offen: Die Bedeutung des Und; die Welt als das, worin beide beisammen sind.
\end{itemize}
